\documentclass[11pt]{article}
\usepackage[margin=1in]{geometry}
\usepackage{amsmath, amssymb}
\usepackage{booktabs}
\usepackage{array}
\usepackage{xcolor}
\usepackage{hyperref}
\usepackage{tikz}
\usetikzlibrary{arrows.meta, positioning, shapes.geometric}
\usepackage{tcolorbox}
\usepackage{enumitem}
\usepackage{parskip}

\hypersetup{
  colorlinks=true,
  linkcolor=blue!50!black,
  urlcolor=blue!50!black,
  citecolor=blue!50!black,
}

\newcommand{\PAF}{\text{PAF}}
\newcommand{\TMREL}{\text{TMREL}}
\newcommand{\SSPtwo}{\text{SSP2{-}RCP4.5}}
\newcommand{\mIHME}{m^{\text{IHME}}}
\newcommand{\mAttr}{D^{\text{attr}}}

\title{Methodology walkthrough \\
\large World Bank climate-health burden projections, 2022--2050}
\author{Aaron Osgood-Zimmerman, World Bank \\ \texttt{aaron.oz@worldbank.org}}
\date{June 2026}

\begin{document}
\maketitle

\begin{abstract}
\noindent
This document explains the math the pipeline performs to project
temperature-attributable mortality and years of life lost (YLLs) under
SSP-RCP climate scenarios out to 2050. Three audiences are addressed:
climatologists running the pipeline (Caspar), epidemiologists checking the
attribution framework (Samuel), and project leads / future analysts
reviewing the deliverable methodology (Charlie and others).

Part~I covers background and formal math: the Burkart 2021 exposure-response
framework, the IHME GBD 2021 forecast scaffolding, and the
\emph{unscaling-rescaling} ratio that translates IHME's reference-scenario
mortality into burden estimates under arbitrary SSP-RCP scenarios. Part~II
walks through concrete worked examples drawn from the pipeline's trace
output (Colombia 2010 Samuel-validated data) and a hypothetical
cross-scenario example, showing each formula applied to real numbers.
\end{abstract}

\tableofcontents

% =====================================================================
\section{Where this document fits}
% =====================================================================

This is the standalone methodology reference for the deliverable. It is
self-contained: a reader following along should not need to look at the
companion documents.

For deeper reading, the supporting docs live in the project repository
at \url{https://github.com/aaron-oz/2026-climate-health-burden-projections}:
\begin{description}[style=multiline,leftmargin=4.2cm]
  \item[\texttt{gbd/ihme-plan-b-prep.tex}] Original methodology memo
    (denser, more abbreviated than this document).
  \item[\texttt{colombia-validation-state.tex}] Empirical validation of
    the pipeline against Samuel Cuervo's published Colombia 2010--2019
    numbers; 18 of 19 headline metrics within $\pm 5\%$.
  \item[\texttt{pipeline-math-verification.md}] Row-by-row independent
    re-derivation of the per-pixel math.
  \item[\texttt{pipeline-state-for-cckp.md}] Pipeline architecture,
    script ordering, input file formats.
  \item[\texttt{caspar-workflow.md}] Step-by-step production run
    instructions.
\end{description}

\begin{tcolorbox}[colback=blue!5,colframe=blue!50!black,title=Scope: what
this document is NOT]
\textbf{Not in scope.} This document does not re-derive Burkart's
MR-BRT exposure-response curve fitting, does not re-derive the IHME GBD
forecast model, and does not re-justify the choice of TMREL definition or
ERF cause set. Those are external inputs documented in Burkart et~al.\
(2021) and the GBD 2021 forecast appendix; we use their published outputs
as given. What \emph{is} in scope is the math we layer on top of those
outputs to produce scenario-specific attributable-burden projections, and
how each step is implemented in the pipeline.
\end{tcolorbox}

% =====================================================================
\section{Notation}
% =====================================================================

\renewcommand{\arraystretch}{1.15}
\begin{center}
\begin{tabular}{l l}
\toprule
Symbol & Meaning \\
\midrule
$c$ & GBD cause (one of 17 temperature-sensitive causes) \\
$l$ & Location (country, GBD location ID) \\
$t$ & Year (1990--2050) \\
$a$ & Age group (GBD 5-year bins: 0--4, 5--9, \ldots, 80+) \\
$s$ & Sex (male, female) \\
$d$ & Draw index ($0, \ldots, N_{\text{draws}}-1$; we use $N=500$) \\
$X$ & Scenario (SSP1-RCP1.9, SSP2-RCP4.5, SSP3-RCP7.0, SSP5-RCP8.5) \\
$\SSPtwo$ & Reference scenario for IHME's published forecasts \\
\midrule
$T_{l,t,p}$ & Daily-mean temperature ($^{\circ}$C) at pixel $p$ on day $t$ \\
$z_{l,p,t}$ & Climate zone at pixel $p$ in year $t$ = round(annual mean $T$) \\
$\tau_{z,t}$ & TMREL (theoretical minimum risk exposure level) for zone
$z$, year $t$ \\
$\pi_{l,t,p}$ & Within-location population weight; $\sum_{p} \pi = 1$
within $(l, t)$ \\
\midrule
$\rho_c(z, T)$ & Raw relative risk for cause $c$, zone $z$, daily
temperature $T$ (Burkart ERF) \\
$\tilde\rho_c(z, T, t)$ & Rescaled RR: $\rho_c(z, T) /
\rho_c(z, \tau_{z,t})$, so $\tilde\rho(z, \tau, t) = 1$ \\
$\PAF^{\text{heat}}_{c,l,t}$ & Heat-attributable PAF for cause $c$,
location $l$, year $t$ \\
$\PAF^{\text{cold}}_{c,l,t}$ & Cold-attributable PAF, same indexing \\
$\tilde S_{c,X,l,t}$ & Pipeline's pop-weighted PAF aggregate under
scenario $X$ ($= \PAF^{\text{heat}} + \PAF^{\text{cold}}$) \\
\midrule
$\mIHME_{c,l,a,s,t,d}$ & IHME's forecast count of cause-$c$ deaths under
their reference \\
$\mAttr_{c,X,l,a,s,t,d}$ & Pipeline's projection: temperature-attributable
deaths under scenario $X$ \\
$\text{ex}_{l,a,s,t}$ & Remaining life expectancy at age of death (from
life tables) \\
\bottomrule
\end{tabular}
\end{center}

The pipeline's intermediate files store these in the same shape:
pixel-day rows for $T$ and $\pi$, zone-keyed rows for $\tau$,
$(z, T, c, d)$-keyed rows for $\rho$, and $(l, a, s, t, c, d)$-keyed
rows for $\mIHME$ and $\mAttr$.

% =====================================================================
\part{The math}
% =====================================================================

\section{Background: how temperature attribution works}

\subsection{Why ``attribution''}

In epidemiology, an \emph{attributable fraction} answers the question
``what share of deaths from cause $c$ would not have occurred if exposure
were at some counterfactual level?'' For non-optimal temperature, the
counterfactual is the TMREL --- the daily temperature at which mortality
risk is minimized for that climate zone. A pixel-day where actual
temperature equals the TMREL contributes zero attributable burden; a
pixel-day at the cold or heat tail contributes positive attributable
burden weighted by how much it raises relative risk.

For an individual pixel-day with population weight $\pi$ and relative
risk $\rho$, the attributable fraction is:
\begin{equation}
\text{AF}_{\text{pixel-day}} = \pi \cdot \frac{\rho - 1}{\rho}
\quad \text{(if } \rho \geq 1)
\label{eq:af-pixel}
\end{equation}
For $\rho < 1$ (a temperature \emph{below} the risk-minimum within a zone,
which can happen for some causes), the sign flips: that pixel-day
contributes a small \emph{negative} attributable burden (it would have
been protective compared to TMREL).

Summing Eq.~\eqref{eq:af-pixel} over all pixel-days in a location-year
that fall on the cold side of TMREL gives the location-year
cold-attributable PAF; summing over heat-side pixel-days gives the
heat-attributable PAF. Multiplying by total cause-specific deaths yields
attributable deaths.

\subsection{Burkart's temperature exposure-response functions}

Burkart et al.~(\textit{Lancet}, 2021) fit relative-risk curves
$\rho_c(z, T)$ for 17 temperature-sensitive causes using meta-regression
(MR-BRT) on cause-specific mortality data from 9 countries (64.9M
deaths). The fit produces 1,000 posterior draws of $\rho_c$ for each
combination of climate zone $z \in [6, 28]\,^{\circ}$C and daily
temperature $T \in [-22, 47]\,^{\circ}$C. We use these curves directly;
the pipeline does not refit them.

The 17 causes split into three broad groups by mechanism:
\begin{itemize}[itemsep=2pt]
  \item \textbf{Cardiovascular} (5): ischemic heart disease, stroke,
    hypertensive heart disease, cardiomyopathy, chronic kidney disease.
  \item \textbf{Respiratory and metabolic} (3): lower respiratory
    infections, chronic obstructive pulmonary disease, diabetes.
  \item \textbf{Injuries} (9): homicide, suicide, mechanical injuries,
    transport-road, transport-other, animal-related, disaster-related,
    drowning, other-unintentional.
\end{itemize}

\subsection{Climate zones and TMREL}

A pixel's climate zone $z$ is the rounded annual mean temperature for
that pixel-year. Zones are integers spanning 6--28$\,^{\circ}$C; the
Burkart ERF curves are zone-specific because the exposure-response
relationship is shaped by acclimatization (people in cold zones tolerate
cold differently than people in hot zones).

The TMREL is the temperature that minimizes \emph{death-weighted}
relative risk for a given zone and location. It is provided by IHME at
1,034 GBD locations and three time points (1990, 2010, 2020); our
pipeline averages these to one zone-specific TMREL used across the
study period.\footnote{For Workflow~B's projections, we use the 2020
TMREL throughout the 2022--2050 horizon. Decadal TMREL projections
beyond 2020 are an open methodology choice; the IHME GBD 2021 forecast
appendix Section~2.1.6.7 discusses a 10-year rolling mean approach.
Within our pipeline the cross-scenario ratio cancels first-order TMREL
drift (since both numerator and denominator use the same TMREL).}

\subsection{The unscale/rescale problem}

IHME's GBD 2021 forecast publishes cause-specific mortality rates under
their reference scenario, which roughly corresponds to $\SSPtwo$. For 12
of our 17 causes, IHME applied a temperature-scalar adjustment to the
forecast (so the published count already reflects $\SSPtwo$ temperature
effects). For the other 5, no temperature adjustment was applied.

We want to project mortality \emph{under other climate scenarios}
(SSP1-RCP1.9, SSP3-RCP7.0, SSP5-RCP8.5). IHME does not publish per-SSP
forecasts. Our task: take the published $\SSPtwo$-anchored mortality and
adjust it to reflect a different climate, using \emph{our} pipeline
applied to bias-corrected CMIP6 temperature for the target scenario.

This is Workflow~B (next section).

% =====================================================================
\section{The pipeline math, step by step}
% =====================================================================

For one (location, year, cause) tuple, the pipeline produces a
$\PAF^{\text{heat}}$ and $\PAF^{\text{cold}}$ value, then multiplies by
cause-specific mortality counts to get attributable deaths.
Per-pixel-day:

\subsection{Climate zone assignment}

For each pixel $p$ in location $l$ in year $t$:
\begin{equation}
z_{l,p,t} = \text{round}\!\left( \frac{1}{N_{\text{days}}} \sum_{\text{day} \in t} T_{l,t,p,\text{day}} \right)
\quad \text{clipped to } [6, 28]
\label{eq:zone}
\end{equation}
The pipeline computes one zone per pixel per year. As climate warms, a
pixel's zone shifts upward over time; this is automatic.

\subsection{TMREL lookup and heat/cold classification}

The TMREL is keyed by (zone, year). For pixel-day $i$ with zone
$z_i$ and daily temperature $T_i$:
\begin{equation}
\text{risk}_i = \begin{cases}
\text{cold} & T_i < \tau_{z_i, t} \\
\text{heat} & T_i > \tau_{z_i, t} \\
\text{at-TMREL} & T_i = \tau_{z_i, t}\ \text{(no contribution)}
\end{cases}
\end{equation}

\subsection{RR rescaling at the TMREL}

This step is easy to overlook. The pipeline divides the raw ERF curve
$\rho_c(z, T)$ by its own value at $\tau_{z,t}$:
\begin{equation}
\tilde\rho_c(z, T, t) = \frac{\rho_c(z, T)}{\rho_c(z, \tau_{z,t})}
\label{eq:rescale}
\end{equation}
After rescaling, $\tilde\rho_c(z, \tau, t) = 1$ exactly. This is the
Burkart 2021 convention: relative risk is reported with the TMREL as the
unit reference, not the curve fit's arbitrary baseline. \emph{All
subsequent PAF computations use $\tilde\rho$, not $\rho$.}

\subsection{Within-location PAF aggregation}

For location $l$, year $t$, cause $c$:
\begin{align}
\PAF^{\text{cold}}_{c,l,t} &= \sum_{i \in \text{cold pixel-days}}
    \pi_{l,t,p_i} \cdot \frac{\tilde\rho_c(z_i, T_i, t) - 1}{\tilde\rho_c(z_i, T_i, t)} \label{eq:paf-cold} \\
\PAF^{\text{heat}}_{c,l,t} &= \sum_{i \in \text{heat pixel-days}}
    \pi_{l,t,p_i} \cdot \frac{\tilde\rho_c(z_i, T_i, t) - 1}{\tilde\rho_c(z_i, T_i, t)} \label{eq:paf-heat}
\end{align}
The pop weights $\pi$ are computed in
\texttt{03\_load\_temperature.R} as
$\pi_{l,t,p,\text{day}} = \text{pop}_{l,t,p} \,/\, \sum_{p,\text{day}}
\text{pop}$, summing to 1 over (pixel $\times$ day) within $(l, t)$. So
the PAF is a pixel-day-weighted average attributable fraction.

The pop-weighted aggregate non-optimal PAF that drives Workflow~B's ratio
is:
\begin{equation}
\tilde S_{c,X,l,t} = \PAF^{\text{heat}}_{c,l,t} + \PAF^{\text{cold}}_{c,l,t}
\quad \text{computed using scenario-}X\text{ temperature data}
\label{eq:S}
\end{equation}

\subsection{Attributable deaths}

For cause $c$, age $a$, sex $s$:
\begin{align}
\mAttr_{c,l,a,s,t}^{\text{cold}} &= m_{c,l,a,s,t} \cdot \PAF^{\text{cold}}_{c,l,t} \label{eq:attrib-cold} \\
\mAttr_{c,l,a,s,t}^{\text{heat}} &= m_{c,l,a,s,t} \cdot \PAF^{\text{heat}}_{c,l,t} \label{eq:attrib-heat}
\end{align}
where $m_{c,l,a,s,t}$ is the cause-specific death count (which the
pipeline reads from its mortality input). The PAF, being a location-year
quantity, broadcasts across all (age, sex) cells. Attributable YLLs are
then $\mAttr \times \text{ex}$.

% =====================================================================
\section{Workflow B: unscaling and rescaling}
% =====================================================================

\subsection{The setup}

We have:
\begin{itemize}[itemsep=2pt]
  \item IHME's published mortality forecast $\mIHME_{c,l,a,s,t}$ ---
    counts derived from rate $\times$ population, at the country level,
    annual, 500 draws each.
  \item Our pipeline applied to the bias-corrected CMIP6 temperature for
    scenario $X$, producing $\PAF^{\text{heat}}_{c,X,l,t}$,
    $\PAF^{\text{cold}}_{c,X,l,t}$, and $\tilde S_{c,X,l,t}$.
  \item Our pipeline applied to $\SSPtwo$ CMIP6 temperature, producing
    the same quantities under the reference scenario.
\end{itemize}

We want $\mAttr_{c,X,l,a,s,t}$, the temperature-attributable deaths from
cause $c$ under scenario $X$.

The challenge: IHME's $\mIHME_{c,l,a,s,t}$ already has $\SSPtwo$-climate
effects baked in (for 12 of the 17 causes). If we just multiply by
$\PAF_X$, we're double-counting some of the temperature signal. We need
to first \emph{unscale} IHME's count to remove the $\SSPtwo$ temperature
contribution, then \emph{rescale} to add scenario-$X$'s contribution.

\subsection{The ratio formula}

The ratio of pipeline-computed pop-weighted PAFs between scenarios is
the scale factor:
\begin{equation}
r_{c,X,l,t} = \frac{\tilde S_{c,X,l,t}}{\tilde S_{c,\SSPtwo,l,t}}
\label{eq:ratio}
\end{equation}
Then for the 12 temp-scaled causes:
\begin{equation}
\mAttr_{c,X,l,a,s,t}^{\text{heat}} = \mIHME_{c,l,a,s,t} \cdot r_{c,X,l,t} \cdot \PAF^{\text{heat}}_{c,X,l,t}
\label{eq:wb-12}
\end{equation}
and analogously for cold. The product $\mIHME \cdot r$ is the
scenario-$X$ ``total mortality envelope'' --- IHME's reference total
adjusted upward (if scenario $X$ has more temperature exposure than
$\SSPtwo$) or downward. Multiplying by scenario-$X$'s heat PAF gives the
heat-attributable share of that adjusted total.

\subsection{Why the ratio cancels calibration drift}

A skeptical reader might object: ``You're using your own pipeline's PAFs
in both numerator and denominator. Don't they need to match IHME's
absolute calibration?''

The answer is no, and it's the key advantage of the ratio formulation.
Suppose our pipeline and IHME's compute PAFs that differ by some
unknown multiplicative factor $k_{c,l,t}$ (because we use Burkart's ERFs
and they use a slightly different fit, or because our bias-corrected
CMIP6 differs subtly from theirs):
\begin{equation*}
\tilde S^{\text{ours}}_{c,X,l,t} = k_{c,l,t} \cdot \tilde S^{\text{true}}_{c,X,l,t}
\quad\text{and}\quad
\tilde S^{\text{ours}}_{c,\SSPtwo,l,t} = k_{c,l,t} \cdot \tilde S^{\text{true}}_{c,\SSPtwo,l,t}
\end{equation*}
Then in the ratio, $k$ cancels:
\begin{equation*}
\frac{\tilde S^{\text{ours}}_{c,X,l,t}}{\tilde S^{\text{ours}}_{c,\SSPtwo,l,t}}
= \frac{\tilde S^{\text{true}}_{c,X,l,t}}{\tilde S^{\text{true}}_{c,\SSPtwo,l,t}}
\end{equation*}
The ratio carries pure cross-scenario information: how much more (or
less) temperature-attributable burden there is under $X$ than under
$\SSPtwo$, in proportional terms. The absolute level is anchored by
IHME's $\mIHME$. This is a stronger position than trying to match IHME's
absolute calibration.

\subsection{The 12 temp-scaled vs 5 non-temp-scaled causes}

The GBD 2021 forecast appendix (Section~2.1.6.7) lists 12 causes for
which IHME applied a temperature-scalar adjustment internally:
\begin{quote}
``ischaemic heart disease, stroke, hypertensive heart disease, diabetes,
chronic kidney disease, lower respiratory infection, chronic obstructive
pulmonary disease, homicide, suicide, mechanical injuries,
transport-related injuries, and drowning.''
\end{quote}
For these 12, Eq.~\eqref{eq:wb-12} applies as stated.

The other 5 causes IHME forecasts (cardiomyopathy and myocarditis,
animal-related injuries, disaster-related injuries, other unintentional
injuries, other transport injuries\footnote{We split IHME's
``transport-related injuries'' umbrella into road and other; both are
treated as temp-scaled.}) are forecast \emph{without} a temperature
adjustment. For these, IHME's $\mIHME$ has no temperature signal to
``unscale,'' so the ratio collapses to identity:
\begin{equation}
\mAttr_{c \in \text{non-temp}, X, l, a, s, t}^{\text{heat}}
= \mIHME_{c,l,a,s,t} \cdot \PAF^{\text{heat}}_{c,X,l,t}
\label{eq:wb-5}
\end{equation}
This is Approach~1 (the simpler textbook attributable-fraction formula)
applied to those 5 causes. From a code perspective it is the same
pipeline; only the post-processing step in
\texttt{util\_workflow\_b\_ratio.R} differs by branching on whether
$c$ is in the temp-scaled set.

\subsection{Workflow diagram}

\begin{center}
\begin{tikzpicture}[
  node distance=1.0cm and 0.8cm,
  box/.style={rectangle, draw, rounded corners, minimum height=0.8cm,
              text width=4cm, align=center, font=\small},
  bigbox/.style={rectangle, draw, rounded corners, minimum height=1.0cm,
              text width=5cm, align=center, font=\small},
  arrow/.style={-Stealth, thick}
]
\node[box] (ihmerate) {IHME forecast rate \\ $r_{c,l,a,s,t}$};
\node[box, right=of ihmerate] (ihmepop) {IHME forecast pop \\ pop$_{c,l,a,s,t}$};
\node[box, below=of ihmerate, xshift=2cm] (mihme) {$\mIHME_{c,l,a,s,t}$ \\ (under their $\SSPtwo$ reference)};
\draw[arrow] (ihmerate) -- (mihme);
\draw[arrow] (ihmepop)  -- (mihme);

\node[box, below=of mihme, xshift=-5cm] (cckp2) {CCKP CMIP6 \\ $\SSPtwo$ temp};
\node[box, below=of mihme, xshift= 5cm] (cckpx) {CCKP CMIP6 \\ scenario-$X$ temp};
\node[bigbox, below=of cckp2] (pipe2) {Pipeline \\ (Eqs.~\ref{eq:zone}--\ref{eq:S}) \\ $\tilde S_{c,\SSPtwo,l,t}$};
\node[bigbox, below=of cckpx] (pipex) {Pipeline \\ (Eqs.~\ref{eq:zone}--\ref{eq:S}) \\ $\tilde S_{c,X,l,t}$, $\PAF^{\text{heat}}_{c,X}$, $\PAF^{\text{cold}}_{c,X}$};
\draw[arrow] (cckp2) -- (pipe2);
\draw[arrow] (cckpx) -- (pipex);

\node[bigbox, below=of mihme, yshift=-4.2cm] (ratio) {Ratio (Eq.~\ref{eq:ratio}) \\ $r_{c,X,l,t} = \tilde S_X / \tilde S_{\SSPtwo}$};
\draw[arrow] (pipe2) -- (ratio);
\draw[arrow] (pipex) -- (ratio);

\node[bigbox, below=of ratio] (out) {
$\mAttr_{c,X,l,a,s,t}^{\text{heat/cold}}$ \\
$= \mIHME \cdot r \cdot \PAF^{\text{heat/cold}}_{c,X}$ \quad (Eq.~\ref{eq:wb-12}) \\
\textit{or, for non-temp-scaled $c$:} \\
$= \mIHME \cdot \PAF^{\text{heat/cold}}_{c,X}$ \quad (Eq.~\ref{eq:wb-5})
};
\draw[arrow] (mihme) -- (out.north);
\draw[arrow] (ratio) -- (out);
\draw[arrow] (pipex.east) to[bend left=45] node[midway, right] {$\PAF^{\text{heat/cold}}_{c,X}$} (out.east);
\end{tikzpicture}
\end{center}

% =====================================================================
\part{Worked examples}
% =====================================================================

The examples below use real values from the pipeline's trace tool applied
to Colombia 2010 with Samuel's validated input data (the dataset where
the pipeline matches Samuel's published numbers on 18 of 19 headline
metrics, $\pm 5\%$). The cross-scenario example in
Section~\ref{sec:wb-example} uses synthesized numbers because we don't
yet have CCKP-driven production runs; the synthesized values are clearly
flagged and will be updated with real production output when Caspar's
runs complete.

% =====================================================================
\section{Within-pipeline trace: one pixel}
% =====================================================================

We trace pixel 527 in Colombia, a Bogotá D.C. metro pixel (subloc 11),
in 2010. Following the math of Section~3:

\textbf{Raw input} (pixel 527, 2010-01-02, from temperature.rds):
\begin{itemize}[itemsep=0pt]
  \item Daily temperature $T = 13.589\,^{\circ}$C
  \item Pixel population: 5{,}100{,}137
  \item Subloc: 11 (Bogotá D.C.)
\end{itemize}

\textbf{Step 1: daily-temp encoding} (pipeline stores
\texttt{daily\_temp\_10 = round($T \cdot 10$)} to avoid floating-point
issues in lookups):
\[
\text{daily\_temp\_10} = \text{round}(13.589 \cdot 10) = 136
\]

\textbf{Step 2: zone} (Eq.~\ref{eq:zone}):
\[
z_{l, p=527, t=2010} = \text{round}(13.786\,^{\circ}\text{C annual mean}) = 14
\]

\textbf{Step 3: pop weight} ($\pi_{l,t,p}$):
\[
\pi = \frac{5{,}100{,}137}{3{,}188{,}693{,}158\ \text{person-days in subloc 11, 2010}} = 1.599 \times 10^{-3}
\]

\textbf{Step 4: TMREL} for $(z=14, t=2010)$:
$\tau = 235$ (= $23.5\,^{\circ}$C).

\textbf{Step 5: raw RR} for cvd\_ihd at $(z=14, T=136)$:
$\rho_{\text{cvd\_ihd}}(14, 136) = 1.018998$.

\textbf{Step 6: rescale RR} (Eq.~\ref{eq:rescale}):
\[
\tilde\rho_{\text{cvd\_ihd}}(14, 136, 2010)
= \frac{\rho_{\text{cvd\_ihd}}(14, 136)}{\rho_{\text{cvd\_ihd}}(14, 235)}
= \frac{1.018998}{0.944913}
= 1.078403
\]
The reference RR at TMREL (0.944913) is below 1 because the un-rescaled
Burkart curve is fit with an arbitrary baseline; after rescaling, the
curve reads RR = 1 at the TMREL by construction.

\textbf{Step 7: heat/cold classification}: $T = 136 < \tau = 235$, so
this pixel-day contributes to the \textbf{cold} side.

\textbf{Step 8: PAF contribution} (Eq.~\ref{eq:af-pixel} with rescaled
$\tilde\rho$):
\[
\text{AF}_{527,\ \text{2010-01-02},\ \text{cvd\_ihd}}
= \pi \cdot \frac{\tilde\rho - 1}{\tilde\rho}
= 1.599\!\times\!10^{-3} \cdot \frac{1.078403 - 1}{1.078403}
= 1.163 \times 10^{-4}
\]

This is one cold-attributable contribution. The full
$\PAF^{\text{cold}}_{\text{cvd\_ihd}, \text{Colombia}, 2010}$ in
Eq.~\eqref{eq:paf-cold} is the sum of such contributions over all
cold-classified pixel-days in 2010. The pipeline's trace tool reproduces
each value above exactly; row-by-row independent re-derivation matches
to 8 decimal places (see
\texttt{pipeline-math-verification.md}).

% =====================================================================
\section{Within-pipeline trace: contrasting pixels}
% =====================================================================

The same trace tool applied to three contrasting pixels in Colombia 2010
illustrates how the math handles different climate regimes. All values
are for cause cvd\_ihd on the pixel's annual-median representative day:

\begin{center}
\small
\begin{tabular}{r r r r r r r r r}
\toprule
pixel & zone & ann.\ mean & $T$ on rep day & TMREL & risk & raw $\rho$ & rescaled $\tilde\rho$ & AF (Eq.\ \ref{eq:af-pixel}) \\
& & $^{\circ}$C & $^{\circ}$C & ($^{\circ}$C) & & & & \\
\midrule
527 (Bogotá) & 14 & 13.8 & 13.6 & 23.5 & cold & 1.019 & 1.078 & $+1.16\!\times\!10^{-4}$ \\
639 (Andes ft.) & 22 & 22.0 & 22.0 & 25.2 & cold & 1.011 & 1.039 & $+5.42\!\times\!10^{-5}$ \\
28 (Caribbean) & 27 & 27.2 & 27.2 & 25.1 & heat & 0.992 & 0.978 & $-3.76\!\times\!10^{-5}$ \\
\bottomrule
\end{tabular}
\end{center}

What this shows:
\begin{itemize}[itemsep=2pt]
  \item \textbf{Pixel 527 (Bogotá, cool high altitude):} $T = 13.6\,^{\circ}$C
    is far below TMREL = $23.5\,^{\circ}$C, so risk is ``cold.'' The
    rescaled RR is 1.078 (~8\% excess risk vs TMREL). This pixel
    contributes positively to cold-attributable burden.
  \item \textbf{Pixel 639 (Andean foothills, moderate):} $T = 22.0\,^{\circ}$C
    is 3 degrees below TMREL = $25.2\,^{\circ}$C; cold but mildly so.
    Smaller positive contribution.
  \item \textbf{Pixel 28 (Caribbean coast, hot tropical):} $T = 27.2\,^{\circ}$C
    is just 2 degrees above TMREL = $25.1\,^{\circ}$C; heat but on the
    flat shoulder of the cvd\_ihd ERF. Rescaled RR is 0.978 ($<1$), so
    the AF is \emph{negative} --- this pixel-day would actually have
    been associated with slightly lower mortality than the TMREL
    reference for cvd\_ihd. The pipeline sums these signed contributions
    in Eqs.~\eqref{eq:paf-cold}--\eqref{eq:paf-heat}, so negative
    contributions partially cancel positive ones.
\end{itemize}

These illustrate why the country-aggregate $\PAF^{\text{heat}}$ for
cvd\_ihd in Colombia 2010 is slightly negative ($-0.22\%$, see next
section): Colombia's heat-side pixel-days are mostly on the flat
shoulder where rescaled RR is close to or below 1.

% =====================================================================
\section{Country aggregate: Colombia 2010 cvd\_ihd}
% =====================================================================

Summing per-pixel-day contributions over all pixel-days in Colombia 2010
(1{,}494 pixels $\times$ 365 days = 545{,}310 rows in the temperature
intermediate), then grouping by (subloc, risk, cause), the pipeline
saves at the country level:

\begin{center}
\renewcommand{\arraystretch}{1.1}
\begin{tabular}{l r r r}
\toprule
& heat & cold & non-optimal \\
\midrule
$\PAF^{\cdot}_{\text{cvd\_ihd, Colombia, 2010}}$ & $-0.22\%$ & $+4.48\%$ & $+4.26\%$ \\
\midrule
& deaths\_heat & deaths\_cold & deaths\_non-opt \\
$\mAttr$ for cvd\_ihd = $\mIHME \cdot \PAF$ & $-65$ & $+1{,}316$ & $+1{,}251$ \\
\bottomrule
\end{tabular}
\end{center}

Out of 29{,}364 total cvd\_ihd deaths in Colombia 2010, the pipeline
attributes 4.3\% (1{,}251 deaths) to non-optimal temperature exposure ---
nearly all on the cold side. The negative heat number reflects the flat
shoulder of the cvd\_ihd ERF on the heat side combined with Colombia's
modest heat-exposure in 2010.

Independent re-derivation: summing my hand-computed AF contributions
across all subloc-11 pixel-days for cvd\_ihd 2010 gives
$\PAF^{\text{cold}}_{\text{subloc 11}} = 0.06956161$, which matches the
pipeline's saved value to 8 decimal places. Multiplying by Colombia's
per-subloc cvd\_ihd death counts and summing yields the country total of
1{,}316 cold-attributable deaths, matching the pipeline's
\texttt{burden\_125.rds} exactly.

% =====================================================================
\section{Cross-scenario worked example (Workflow B ratio)}
\label{sec:wb-example}
% =====================================================================

\textit{Caveat: this section uses \textbf{synthesized illustrative
values} for the cross-scenario numbers. Real production values will
substitute these once Caspar's runs complete; the math is unchanged.}

\subsection{Setup}

Take Colombia in year 2030, cause cvd\_ihd, all ages \& sexes. Suppose
IHME publishes a forecast count of cvd\_ihd deaths of:
\[
\mIHME_{\text{cvd\_ihd, Colombia, 2030}} = 38{,}000 \text{ deaths}
\quad \text{(synthesized; IHME's actual forecast TBD)}
\]
This is under their reference scenario, approximately $\SSPtwo$.

Run our pipeline twice: once with CCKP CMIP6 $\SSPtwo$ daily temperature
for Colombia 2030, once with CCKP CMIP6 SSP5-RCP8.5 daily temperature for
Colombia 2030. Suppose the pipeline outputs:

\begin{center}
\renewcommand{\arraystretch}{1.1}
\begin{tabular}{l r r}
\toprule
& Reference $\SSPtwo$ & Target SSP5-RCP8.5 \\
\midrule
$\PAF^{\text{heat}}_{\text{cvd\_ihd, Colombia, 2030}}$ & $+0.5\%$ & $+1.8\%$ \\
$\PAF^{\text{cold}}_{\text{cvd\_ihd, Colombia, 2030}}$ & $+4.2\%$ & $+3.2\%$ \\
$\tilde S$ (heat + cold)                           & $+4.7\%$ & $+5.0\%$ \\
\bottomrule
\end{tabular}
\end{center}

In this hypothetical: SSP5-RCP8.5 in 2030 has more heat-attributable
burden (more hot days above TMREL) and slightly less cold-attributable
burden (fewer cold days below TMREL). The non-optimal total is just
slightly higher under SSP5 than $\SSPtwo$ --- the heat increase mostly
cancels with the cold decrease.

\subsection{Apply the ratio}

From Eq.~\eqref{eq:ratio}:
\[
r_{\text{cvd\_ihd, SSP5, Colombia, 2030}}
= \frac{\tilde S_{\text{SSP5}}}{\tilde S_{\SSPtwo}}
= \frac{0.050}{0.047}
= 1.0638
\]
The total mortality envelope grows ~6.4\% under SSP5 vs $\SSPtwo$
because the temperature signal is slightly stronger overall. cvd\_ihd
is one of the 12 temp-scaled causes, so Eq.~\eqref{eq:wb-12} applies:
\begin{align*}
\mAttr_{\text{cvd\_ihd, SSP5, 2030}}^{\text{heat}}
&= \mIHME \cdot r \cdot \PAF^{\text{heat}}_{\text{SSP5}} \\
&= 38{,}000 \cdot 1.0638 \cdot 0.018 \\
&= 727 \text{ heat-attributable deaths} \\[2pt]
\mAttr_{\text{cvd\_ihd, SSP5, 2030}}^{\text{cold}}
&= \mIHME \cdot r \cdot \PAF^{\text{cold}}_{\text{SSP5}} \\
&= 38{,}000 \cdot 1.0638 \cdot 0.032 \\
&= 1{,}294 \text{ cold-attributable deaths} \\[2pt]
\mAttr_{\text{cvd\_ihd, SSP5, 2030}}^{\text{non-opt}}
&= 727 + 1{,}294 = 2{,}021 \text{ non-optimal-temperature deaths}
\end{align*}

\subsection{Compare to $\SSPtwo$ baseline}

Under $\SSPtwo$ at the same location-year, applying the same formula
(with $r = 1$ since reference equals reference):
\begin{align*}
\mAttr_{\text{cvd\_ihd}, \SSPtwo, 2030}^{\text{non-opt}}
&= 38{,}000 \cdot 1.0 \cdot 0.047
= 1{,}786 \text{ deaths}
\end{align*}

So SSP5-RCP8.5 produces ~235 more cvd\_ihd deaths attributable to
non-optimal temperature in Colombia 2030 than $\SSPtwo$ does, with the
\emph{composition} shifting toward heat (heat goes from ~$190$ to ~$727$,
cold goes from ~$1{,}596$ to ~$1{,}294$).

\subsection{The same exercise for a non-temp-scaled cause}

For cardiomyopathy (cvd\_cmp), which is one of the 5 non-temp-scaled
causes, IHME's forecast has no implicit temperature signal. The ratio
collapses to identity per Eq.~\eqref{eq:wb-5}. Suppose
$\mIHME_{\text{cvd\_cmp, Colombia, 2030}} = 5{,}000$ and our pipeline
gives $\PAF^{\text{heat}}_{\text{cvd\_cmp, SSP5}} = 0.4\%$,
$\PAF^{\text{cold}}_{\text{cvd\_cmp, SSP5}} = 1.8\%$ (synthesized
values). Then:
\begin{align*}
\mAttr_{\text{cvd\_cmp, SSP5, 2030}}^{\text{heat}}
&= 5{,}000 \cdot 0.004 = 20 \text{ deaths} \\
\mAttr_{\text{cvd\_cmp, SSP5, 2030}}^{\text{cold}}
&= 5{,}000 \cdot 0.018 = 90 \text{ deaths}
\end{align*}
No ratio applied; the formula reduces to ``IHME's mortality $\times$
scenario-$X$ PAF.''

% =====================================================================
\section{Validation evidence}
% =====================================================================

The math above is verified empirically by reproducing Samuel Cuervo's
published Colombia 2010--2019 temperature-attributable-mortality numbers
to within $\pm 5\%$ on 18 of 19 headline metrics (\texttt{colombia-validation-state.tex} in the project repo). The one remaining discrepancy
(road-injury heat YLLs, $-11.7\%$) is traced to upstream choices in
Samuel's input data that we do not replicate; headline totals and top-3
cause rankings all match.

Row-by-row independent re-derivation of the per-pixel math against the
pipeline's own saved values (Section~7 above; full details in
\texttt{pipeline-math-verification.md}) confirms every formula in
Section~3 is correctly implemented: 9 calculation stages match to 8+
decimal places, and the within-subloc aggregation reproduces the
pipeline's saved per-subloc PAF (0.06956161 vs 0.06956161 for Bogotá D.C.
cvd\_ihd cold, 2010).

% =====================================================================
\section{Reproducing}
% =====================================================================

The pipeline runs from \texttt{global-scripts/run\_location.R} for a
single location, or \texttt{global-scripts/util\_run\_global.R} for a
production grid across locations. The Workflow~B ratio applicator is
\texttt{global-scripts/util\_workflow\_b\_ratio.R}. The trace tool that
generated the values in Sections~7 and~8 is
\texttt{global-scripts/util\_trace\_pipeline.R}.

Operational instructions for the full production run, including the
single-combo benchmark and parallelization pattern, are in
\texttt{caspar-workflow.md} in the project repository.

\bigskip
\hrule
\smallskip

\textit{Source repository:}
\url{https://github.com/aaron-oz/2026-climate-health-burden-projections}

\textit{Key references:}
Burkart et al., ``Estimating the cause-specific relative risks of
non-optimal temperature on daily mortality,''
\textit{The Lancet}, 2021;
Institute for Health Metrics and Evaluation, ``GBD 2021 Forecasting
Methods Appendix,'' Section 2.1.6.7;
World Bank Climate Change Knowledge Portal, CMIP6 daily-x0.25 data
product.

\end{document}
